\section{Dimensionality reduction}

When dealing with big data (also not very big) we'll always clash with the
\textbf{Curse of dimensionality}:
\begin{center}
    \textit{In applied mathematics, curse of dimensionality refers to the problem 
    caused by the exponential increase in volume associated with adding extra 
    dimensions to a mathematical space.}
\end{center}
\vspace{0.7cm}

Let's look at the \textbf{example} of categorizing some bacteria:
\begin{figure}[h]
    \centering
    \includegraphics[width=0.3\textwidth]{dr/example_cd_1d.png}
\end{figure}

To separate the groups better we add new information about another measure
\textcolor{gray}{(before we have that each point has one feature, while now 
we're considering two features for each point, if you think about it you are 
in fact adding information and making the points more different from each other)}
\begin{figure}[h]
    \centering
    \begin{subfigure}[t]{0.2\textwidth}
        \centering
        \includegraphics[width=0.95\linewidth]{dr/example_cd_2d.png}
    \end{subfigure}
    \hspace{0.5cm}
    \begin{subfigure}[t]{0.2\textwidth}
        \centering
        \includegraphics[width=0.95\linewidth]{dr/example_cd_2d_expl.png}
        {
            \captionsetup{width=1.5\linewidth, margin={0cm, -3cm}, labelformat=empty}
            \caption{If we try to maintain for the 2-d scenario the same density of 
            the 1-dimension case, we would end up with 81 samples.}
        }
    \end{subfigure}
\end{figure}

If we extend the space up to the 3th dimension we are capable of 
distinguish the data using only a linear model. To do the same in 2 dimensions
we should have used more complex models.
\begin{figure}[h]
    \centering
    \begin{subfigure}[t]{0.3\textwidth}
        \centering
        \includegraphics[width=0.9\linewidth]{dr/example_cd_3d.png}
    \end{subfigure}
    \hspace{0.2cm}
    \begin{subfigure}[t]{0.3\textwidth}
        \centering
        \includegraphics[width=0.9\linewidth]{dr/example_3d_divison.png}
    \end{subfigure}    
\end{figure}

\newpage
Another phenomena of the curse of dimensionality is \textbf{Distance concentration}: as 
the number of samples goes to infinity, the relative distances between point tend to 
become zero, starting to concetrate on the expected value.
\begin{figure}[h]
    \centering
    \includegraphics[width=0.7\textwidth]{dr/distance_concentration.png}
\end{figure}

\vspace{2cm}
\subsection{Principal Component Analysis (PCA)}

What if we have $n$ data points, $p$ features and $n<p$ ?
It means that some fatures might look something completely different, but in 
reality some of them are just a combination of others $\rightarrow$ we 
have \textbf{collinear features}.
Collinear features make harder our job of understanding how each feature influences
the targets : let's assume that $x_1$ and $x_2$ are collinear
\begin{equation}
    y = \beta_0 + \beta_1 x_1 + \beta_2 x_2 +\varepsilon
\end{equation}
due collinearity, we can't see the individual effect of $x_1$ and $x_2$ on $y$.
%\vspace{0.8cm}

\begin{figure}[h]
    \centering
    \includegraphics[width=0.5\textwidth]{dr/pca/dr_matrix_example.png}    
\end{figure}
What we want to do is to summarize the data having $p$ features with the smaller
set ok $k$ artifically derived features; to do this we must balance the clarity of
representation and the loss of information due the compression $\Rightarrow$ 
\textbf{Principal Component Analysis (PCA)} aims to do this using linear 
combinations of the selected variables.

\vspace{0.6cm}
\textcolor{Green}{\textbf{Example}:}\\
Imagine that we want to describe our blood cells sample, and the initial 
features are \textit{diameter} and \textit{thickness}, how can we reduce
the number of features?
\begin{figure}[h]
    \centering    
    \includegraphics[width=0.3\textwidth]{dr/pca/blood_cells.png}
\end{figure}
\newpage

We can think of projecting the points along one of the two dimensions, choosing it
based on how spreaded the data are along it (e.g. so preserving more information).
\begin{center} \textcolor{red}{
    \textbf{In statitics information = variability.}   
} \end{center}
\begin{figure} [h]
    \centering
    \begin{subfigure}[c]{0.25\textwidth}
        \centering
        \includegraphics[width=\linewidth]{dr/pca/dr_blood_ex1.png}
    \end{subfigure}
    \hspace{1cm}
    \begin{subfigure}[c]{0.2\textwidth}
        \centering
        \includegraphics[width=\linewidth]{dr/pca/dr_blood_ex2.png}
    \end{subfigure}
\end{figure}

So we need a way to capture the maximal variability of the data: in other
words to find \textit{along which dimension the data spread the most}.
\begin{figure}[h]
    \centering
    \includegraphics[width=0.3\textwidth]{dr/pca/pca_maxvar_example.png}
\end{figure}

PCA's goal is to project the data along those directions of maximal variability,
ending up with a new reference system for the data points \textcolor{gray}{(it's 
nothing different from vector space rotations, in 2D you visualize it by immaging
that the axis rotate)}.\\
\hfill \\

This objective can be expressed in two equivalent ways:
\begin{description}
    \item[Maximize variance]: maximize the data spread along that
    direction
    \item[Minimize residuals]: minimize the distance of the data points
    from the the dimension
\end{description}
\begin{figure}[h]
    \centering
    \includegraphics[width=0.45\textwidth]{dr/pca/pca_goal_visualzed.png}
\end{figure}

In order to do the PCA:
\begin{enumerate}
    \item get a set of data points $x^{(1)}\dots x^{(n)}\in \real^p$
    \item center the data in $(0,0)$ if needed: $x^{(i)}\leftarrow x^{(i)}-
    \bar{x}$
    \item choose a direction $u\in\real^p\ :\ \|x\|_2=1$
    \item project the data $x^{(i)}$ on $u$ : $z_i = u^Tx^{(i)}\in \real$
\end{enumerate}

Taking the objective of maximizing the variance, we observe that since the data are 
centered (avg = 0), the varince is:
\begin{equation*}
    Var(z) = \frac{1}{n}\sum_{i=1}^n z_i^2 = \frac{1}{n}\sum_{i=1}^n 
    (u^Tx^{(i)})^2 = 0
    \ \rightarrow\ 
    \max_{u\in\real^p}\frac{1}{n}\sum_{i=1}^n (u^Tx^{(i)})^2
\end{equation*}
if we define the \textit{empirical covariance matrix} as
\begin{equation*}
    \Sigma = \frac{1}{n}\sum_{i=1}^n x^{(i)}x^{(i)T} \ \in 
    \real^{p\times p}
\end{equation*}
we get that
\begin{equation*}
    \frac{1}{n}\sum_{i=1}^n (u^Tx^{(i)})^2 = u^T(\frac{1}{n}
    \sum_{i=1}^n x^{(i)}x^{(i)T})u = u^T\Sigma u
    \rightarrow
    \max_{u\in\real^p} u^T\Sigma u
\end{equation*}
The constraint on $u$ can be expressed as $u^Tu=1$ (turning constraint into 
an equation)

By switching to the Lagrangian relaxed problem:
\begin{equation*}
    \mathcal{L}(u,\lambda)=u^T\Sigma u-\lambda(u^Tu-1)
\end{equation*}

and taking the derivative
\begin{equation*}
    \nabla_u (u^T\Sigma u) = 2\Sigma u,\ \nabla_u (u^Tu) = 2u
    \Rightarrow \nabla_u \mathcal(L) = 2\Sigma u-2\lambda u = 0
    \Rightarrow \Sigma u = \lambda u
\end{equation*}

$\Rightarrow$ a candidate direction $u$ must be an \textbf{eigenvector of $\mathbf{\Sigma}$}.\\
So the objective function becomes $u^T\Sigma u=u^T (\lambda u) = \lambda
(u^T u) = \lambda \rightarrow$ by choosing the highest $\lambda$ we get the
highest objective function. So as principal component we take the eigenvector with the highest 
eigenvalue. If we want more dimension, we take also the second, the third ... \\
\hfill \\

We can interpret $\Sigma$ so that, when applied to the data
modifies the shape, it strech them toward their principal component, which 
are the $\Sigma$'s eigenvectors.
\begin{figure}[h]
    \centering
    \includegraphics[width=0.3\textwidth]{dr/pca/covmat_projection.png}
\end{figure}

Two nice properties of PCA are:
\begin{enumerate}
    \item PCA results are rotation-invariant (depend only on the data, not 
    on the starting axis)
    \item PCs are independent/uncorrelated from each other by construction (they are eigenvectors)
\end{enumerate}

\newpage
\subsubsection{Express data with PCA's components}
We can see that the original observations can be re-written using the new
components: given the starting dimensions $\langle x_1\dots x_p \rangle$, we 
have
\begin{equation*}
    PC_j = \phi_{j1}x_1 \dots + \phi_{jp}x_p \Rightarrow z_{ij} = \phi_{1j}x_{i1}
    \dots + \phi_{pj}x_{ip}
\end{equation*}
where $z_{i,j}$ is the projection of $x^{(i)}$ on the $PC_j$ called \textbf{score}
and $\phi_{1j},\dots \phi_{pj}$ are the \textbf{weights} given to each original 
feature to translate it into the new PCA coordinates system. 
\begin{figure}[h]
    \centering
    \includegraphics[width=0.7\textwidth]{dr/pca/pca_projection_example.png}
\end{figure}

\hfill \\
Note: standardization ($x-\mu/\sigma$) address the problem of
high-variance features dominating principal components.
\vspace{0.5cm}


\textbf{Why this stuff is useful other than re-shaping the data?}\\
Usually the data variability depends only on a small subset of all $p$ features, 
so we can use these to have a $k$-low-dimensional/compressed representation of 
the dataset. The choice of $k$ is the trade-off.

\subsubsection{Portion of Variance Explained}
We would like to quantify how much each singular Principal Component (PC)
explains the data variability.
\begin{equation*}
    Var_{\scriptscriptstyle TOT}= \sum_{j=1}^p Var(x_j) = \sum_{j=1}^p \left(\frac{1}{n}\sum_{i=1}^n x_{ij}^2\right)
\end{equation*}
If we name $Var_k$ := variance explained by the single k-th PC and $Var_{exp}$ := the
portion of Variance Explained by the selected first $k$ components
\begin{equation*}
    Var_k = \frac{1}{n}\sum_{i=1}^n z_{ik}^2 \Rightarrow Var_{exp} = \cfrac
    {\sum_k Var_k}{Var_{\scriptscriptstyle TOT}}
\end{equation*}
However, since the PCs are the eigenvalues of $\Sigma$:
\begin{equation*}
    Var_k = \frac{1}{n}\sum_{i=1}^n z_{ik}^2 = u_k^T\Sigma u_k = u_k^T(\lambda_k
    u_k) = \lambda_k
\end{equation*}
\begin{equation*}
    Var_{\scriptscriptstyle TOT} = \sum_{j=1}^p Var(x_j) = tr(\Sigma) = \sum_{j=1}^p \lambda_k
\end{equation*}
So we can define:
\begin{description}
    \item[$\mathbf{EVR_k}$ Explained Variance Ratio of $\mathbf{k}$] = $\cfrac{\lambda_k}
    {\sum_{m=1}p}\lambda_m$
    \item[$\mathbf{CumEV(K)}$ Cumulative Explained Variance] of the $K$ selection of
    component = $\cfrac{\sum_{k=1}^K\lambda_k}{\sum_{m=1}p}\ \lambda_m$   
\end{description}

\newpage
Depending on how much variability we'd like to retain, we can module the number
of dimensions in $K$ based on the \textit{CumEV index}.
This technique can be used also to other purposes, not only to tune the variance
level: it can be used also to set the data noise treshold to remove those components which 
"make too much noise".
\begin{figure}[h]
    \centering
    \includegraphics[width=0.4\textwidth]{dr/pca/cev_visualized.png}
\end{figure}


\subsubsection{How to interpret PCs}
The PC's weights can be used to interpret the meaning of the PCs itself: by loking at them we can 
have a clue about which dimensions go along in the same direction, and from them we can give to the PC some kind of interpetation
\begin{figure}[h]
    \centering
    \includegraphics[width=0.4\textwidth]{dr/pca/interpret_pcs.png}
\end{figure}
\hfill\\

Another useful tool to interpret the PCs' meaning is the \textbf{biplot}: using
PCs as axis, we plot both the points projected on the PCA dimensions and the original
features.
\begin{figure}[h]
    \centering
    \includegraphics[width=0.4\textwidth]{dr/pca/biplot_example.png}
\end{figure} \hfill \\
and we interpret the meaning as following:
\begin{itemize}
    \item \textbf{arrow direction}: along which direction the feature increases
    \item \textbf{arrow length}: how strongly the features is represented by
    these PCs
    \item \textbf{angle between arrows}: relationship between original features
\end{itemize}



\subsection{Non-Linear dimensionality reduction}
In real-life data usually we can assume the \textbf{Manifolds assumption}:
\begin{center}
    \textit{High-dimensional real data is not scattered uniformly across the ambient 
    space (i.e., the space spanned by the d-original dimensions). Rather, it often 
    lies on (or near to) a much lower dimensional CURVED manifold.}    
\end{center}
\begin{figure}[h]
    \centering
    \includegraphics[width=0.2\textwidth]{dr/non_linear/manifold_vis.png}
\end{figure}
PCa is a linear DR techinque, but for real data we need techniques able to properly
model complex non-linear shapes.

%-----------------------------------------------------------------------------------------
\subsubsection{Kernel PCA methods}
This family of methods follows the idea that \textit{data might be non-linear in $\real
^D$ but might become "more linear" after a non-linear feature mapping $\phi(\cdot)$}.
The non-linear transformation is called \textbf{bias expansion} $\phi: \real^p \rightarrow 
\real^q$ and can be choosen accordingly to the analyst's assumption about the shape of the feature; 
usually we take $q>p$ and transform the sampled data into an higher-dimensional space.\\
\hfill \\

We define the basis-expanded matrix $\Phi=(\phi(x_1)\dots \phi(x_n))^T\rightarrow
\Sigma=\Phi^T\Phi$.\\ However, for highly-dimensional projections this approach could
be computationally expensive. In those cases we can take advantage of the
\textbf{kernel trick}: we can compute $\Sigma_{i,j}=k(x_i,x_j)$ using a \textit{kernel 
function $k(\cdot)$}, which measures the similarity between two projections and it's 
defined by $k(x_i,x_j)=\langle\phi(x_i),\phi(x_j)\rangle$ \textcolor{gray}{(scalar product)}.\\
\hfill \\

One of the advantages of this method is the ability of untangling circles.
\begin{figure}[h]
    \centering
    \includegraphics[width=0.8\textwidth]{dr/non_linear/nldr_untangling_circles.png}
\end{figure}
\vspace{0.8cm}

The problem with Kernel PCA is the original space reconstruction: with \textit{linear
PCA} the inverse is explicit and it's exact if projection and original space have the
same dimension, approximante if the projection space has less dimensions than the 
original one.\\
In \textit{Kernel PCA} however $\phi$ is learned/induced by the data
and this could end up in having complex inverse transformations, if not even
non-invertible $\phi$s. Sometimes the only solution to have some reverse-mapping is to rely on
learning the pre-image of the transformation, or more commonly an approximation of it.

\begin{figure}[h]
    \begin{subfigure}{0.48\textwidth}
        \includegraphics[width=0.95\linewidth]{dr/non_linear/inv_kernelPCA.png}
        \caption*{The idea of non-linear kernel PCA mapping.}            
    \end{subfigure}
    \begin{subfigure}{0.5\textwidth}
        \includegraphics[width=0.95\linewidth]{dr/non_linear/kernel_reconst.png}
        \caption*{Example of non-biective kernel PCA, with loss in the original feature 
        reconstrucion.}    
    \end{subfigure}
\end{figure}
\newpage
\hfill\\
%-----------------------------------------------------------------------------------------
\vspace{0.5cm}
\subsubsection{t-distributed Stochastic Neighbour Embedding (t-SNE)}
The key idea of this family of methods is to \textit{preserve the data geometry by
focusing on local neighbours (points near to each other should remain close to each other after
the projection)}.
\begin{figure}[h]
    \centering
    \includegraphics[width=0.6\textwidth]{dr/non_linear/global_vs_local_neigbhbours.png}
    \caption{On the left the distance computed accordingly to global distance; on the
    right the distance is computed using local neighbour.}
    \label{fig:exlocdist}
\end{figure}\hfill \\
Note: This type of methods are really useful for data visualization.\\
\hfill \\

As shown in the Fig \ref{fig:exlocdist} the problem with global similarity is that doesn't
strictly follows the data shape, while based on similarity between point we \dots
watch the recording. 

To do that, we don't need to preserve the distances with respect
to allo the points, but just with his neighbours. This allow to preserve the 
geometric relationship of the data. 
\begin{figure}[h]
    \centering
    \includegraphics[width=0.6\textwidth]{dr/non_linear/nlpca_shape_preservation.png}
\end{figure}
\hfill \\

t-SNE is mainly focused on preserving neighbourhood: points that are close in the
original space should stay clore in the mapping, otherwise they aren't forced to 
preserve the "right far distance" (not blobal distance as focus).
\hfill \\
I translate the distance into probability: if two points are close, they have higher
probability of being neighbours.
To determine the neighbourhood concept, the t-SNE firstly determines the \textbf{
similarity score} between points based on the distance between them and a \textit{
probability distribution}, which allows to decrease the probability smoothly base 
on the distance between points
\begin{figure}[h]
    \centering
    \includegraphics[width=0.55\textwidth]{dr/non_linear/tSNE-similarity.png}
\end{figure}

For each pair of points $i,j$ we could define the \textbf{Gaussian neighbourhood 
distribution} (fore example) as \textcolor{gray}{(we could think about is as "the 
probability that $x_i$ would pick $x_j$ as neighbour")}:
\begin{equation}
    p_{j|i}= \begin{cases}
        \cfrac{exp(-\| x_i - x_j \|^2/(2\sigma_i)^2)}{\sum_{j\neq i}exp(-\| x_i - x_j 
        \|^2/(2\sigma_i)^2)} & j\neq i \\
        0 & j=i
    \end{cases}
\end{equation}
then to make the relationship symmetric $p_{ij} = \frac{p_{i|j}+p_{j|i}}{2}$.\\
To project as much as possible these similarities on a low-dimensional spaces, instead
of using normal distribution t-SNE use \textit{t-student} to discern due its heavier
tails

For the low-dimensional space we use instead a \textit{t-student distribution}: by 
projecting the data on fewer dimensions, the points will end up being more squashed 
together, making the Normal distribution not suited (it will end up considering all the 
points in the same cluster). The \textit{t-student} instead having heavier tails make
the probability distribution more evenly spread, leveraging the dimensionality reduction
squash effect (also called \textit{crowding problem}) $\rightarrow$ thanks to the heavier
tails, low-D points can be pushed away without making $q_{ij}$ too small.
\begin{equation}
    q_{ij}= \begin{cases}
        \cfrac{(1-\| y_i - y_j \|^2)^{-1}}{\sum_{j\neq i}(1-\| y_i - y_j \|^2)^{-1}} 
        & j\neq i \\
        0 & j=i
    \end{cases}
\end{equation} 
\hfill \\

Using different distributions and different dimensionality, we end up with two different
\textbf{similarity matricies}, and out objective becomes making them look alike as much 
as possible $\rightarrow$ minimizingthe mismatch between $p_{ij}$ and $q_{ij}$. To do
tath we iterate to end up with a low-D similarity matrix that is similar to the original
one.
\begin{figure}[h]
    \centering
    \includegraphics[width=0.6\textwidth]{dr/non_linear/similmat_comparison.png}
\end{figure}
\begin{equation}
    C(Y) = KL(P | Q) = \sum_{i\ne j} p_{ij}\log \frac{p_{ij}}{q_{ij}}    
\end{equation}
the difference is computed using \textit{KL divergence}, which is a way to measure distance
between probabilities.
\\
This cost heavily penalizes high mistatch and has strong local structure preservation 
\textcolor{gray}{(focus on maintainig true neighbours, not in preserving the true 
far-away distances)} $\rightarrow$ this might cause \textbf{distortion of the global
structure}.
\\
\textbf{Perplexity on hyperparameters}\\
How do we choose hyperparameters? By changing the variance of the distributions we shape
the concept of neighbourhood, changing the preplexity ($perp. ~ \sigma^2$) $\rightarrow$ 
we have to balance between the dimensionality of the data and  how many cluster do we 
expect to have.
\begin{figure}[h]
    \centering
    \includegraphics[width=0.6\textwidth]{dr/non_linear/hyperparam_perplexity.png}
\end{figure}

%-----------------------------------------------------------------------------------------
\subsubsection{Uniform Manifold APproximantion (UMAP)}
The idea of this method is still to preserve neighbourhood but using different approach:
we build a graph in the high-D space linking the points and we try to apply the same graph
on the low-D space to preserve the structure as much as possible.
\begin{figure}[h]
    \centering
    \includegraphics[width=0.5\textwidth]{dr/non_linear/umap_concept.png}
    \caption{High-D and low-D graphs using UMAP}
\end{figure}
To build the graph, we place circles centered on the data points, and those points having
overlapping circles are neighbours $\rightarrow$ the circle's radius is critical
\begin{figure}[h]
    \centering
    \includegraphics[width=0.4\textwidth]{dr/non_linear/umap_radius.png}
\end{figure}
To overcome this problem we can extend the algorithm using a \textit{local radius}: instead
of fixed value, for some point we use the distance from its n-th neighbour as radius, and
using somethig similar to a distance as weight of the edge.
\\
Differently to \textit{t-SNE} and its perplexity, this method tunes two different parameters:
\begin{description}
    \item[Nearest neighbour]: the number of neighbours used for the high-D radius, tunes how
    "wide" my clusters will be
    \item[Min distance]: tunes how much tightly data points are gonna be projected in the 
    low-D space \textcolor{gray}{(I'm putting some external knwoledge to make the data plots
    more clear, very dangerous for inference)}.
\end{description}
\hfill \\
UMAP algorithm is way faster that t-SNE, uses a more efficient to shuffle the data.

%-----------------------------------------------------------------------------------------
\subsubsection{AutoEncoders (AE)}
The problem with \textit{Kernel PCA} is that is really difficult to turn back in the
original data space, while \textit{t-SNE/UMAP} are good for visualization but we can't
do inference with them. The idea is to use a new class fo algorithms capable of going
back and forth between high-D and low-D spaces: \textit{Autoencoders}, composed by two
subfunctions, \textit{encoders} and \textit{decoders}, which learn a non-linear mapping 
between spaces.
\begin{description}
    \item[Encoder (compression)]: \hfill \\
    $z_i=f_\theta(x_i),\ f_\theta:\real^D\rightarrow\real^d$ with $z=f_\theta(x)=\sigma(Wx+b) 
    \rightarrow$ the conders apply a non-linear function \ssigma := activation function on 
    pre-linear-combination $Wx+b$ (\ssigma, $W$, $b$ are parameters)
    \item[Decoders (reconstruction)]: \hfill \\
    $\hat{x}_i=g_\phi(z_1),\ g_\phi:\real^d\rightarrow \real^D$ with $\hat{x}=g_\phi(z)=
    \rho(Vz+c)$; the structure is the same as encoder, just in the other way around (and 
    uses \diff params).
\end{description}
\begin{figure}[h]
    \centering
    \includegraphics[width=0.35\textwidth]{dr/non_linear/autoencoders__sketch.png}
\end{figure}
We can tune the parameters to lose as less information as possible during compression/
reconstruction phase, maybe even remove only the information we want to loose $\rightarrow$
Autoencoders allow us to decide what we lose.\\
\vspace{0.5cm}

\textbf{Training objective}: to tune the parameter sets $\theta,\phi$ in order to make the
reconstruction as losless as possible, we have to \textit{train the algoirthm}:
\begin{equation}
    \hat{x}_i=g_\phi(f_\theta(x_i)) \rightarrow \min_{\theta,\Phi} \frac{1}{n} 
    \mathcal{L}(x_i,\hat{x_i})
\end{equation}
Commonly we choose $\mathcal{L}(x_i,\hat{x}_i)=\| x_i-\hat{x}_i \|_2^2$.

\hfill\\
\textcolor{gray}{
This method doesn't allow to explain the data variabilty induced by each dimension like in 
PCA, however stiil allow to tune the level the information/varibaility we want to keep
$\rightarrow$ very useful to identify data anomaly: if a point is very far away the original
space will be projected far away, and also during the reconstrucion phase the decoder will try
to reconstruct it as normal point, giving high reconstruction error and displaying on which 
dimension the anomalies are. 
}

\subsubsection{How to choose non-linear DR methods}
Based on your goal:
\begin{itemize}
    \item \textbf{Data visualization}: use \textit{UMAP} as default, good balance between 
    execution time and accuracy (sometimes preserves also the global organization); use 
    \textit{t-SNE} when your main concern is displaying local cluster separation.
    \item \textbf{Nonlinear feature engineering (reusable mapping)}: \textit{autoencoders}
    are the best to find explicit mapping; \textit{Kernel PCA} is a good fit for medium-sized
    dataset for which we have some domain model.
    \item \textbf{Data interpretability (like PCA loadings)}: PCA, that's it    
\end{itemize} 
\vspace{0.3cm}
Method recomendations:
\begin{description}
    \item[Kernel PCA]\hfill
        \begin{description}
            \item[\textcolor{green}{Pros}]: linearize non-linear structure if a kernel can and
            the dataset sie is manageable.
            \item[\textcolor{red}{Cons}]: not suited for large datasets, and urge the need for
            inverse function, usually only computable trough pre-image approximation. 
        \end{description}
    \item[t-SNE/UMAP]\hfill   
        \begin{description}
            \item[\textcolor{green}{Pros}]: use them for visual exploration of local clusters.
            \item[\textcolor{red}{Cons}]: due the global structure distrortion, they can't be
            used to draw conclusions and doing inference.
        \end{description}
    \item[AutoEncoders]\hfill
        \begin{description}
            \item[\textcolor{green}{Pros}]: learn mapping between spaces, allowing reconstruction
            or interpretability capabilities.
            \item[\textcolor{red}{Cons}]: complexity overhead, with smaller datasets there might
            be overfitting.
        \end{description}   
\end{description}